%!TEX root = ../thesis-main.tex
\begin{abstract}
    \emph{Discrete Event Simulation} constitutes a fundamental technique for
    analysing the behaviour of complex and dynamic systems. However, the
    architectural efficiency of a simulator is frequently constrained by the
    mechanisms employed to manage the dependencies among the simulated
    entities. Traditional approaches typically rely on state polling through
    the maintenance of explicit dependency graphs. These methods, while
    performant, often result in monolithic engines where the logic for tracking
    interactions could potentially obscure simulation dynamics, leading to
    fragility in highly dynamic scenarios where the network topology changes
    frequently.

    This thesis addresses these challenges by proposing a reactive architecture
    for the \emph{Alchemist} simulator, through the analysis, the design and
    implementation of a custom, reactive framework tailored to the strict
    \emph{causality constraint} of discrete-event, stochastic simulation. By
    making the simulation environment observable, the responsibility of
    dependency tracking is shifted from the central engine to reactions
    themselves, effectively delegating to reactions themselves the job of
    rescheduling when updates from their dependencies are detected.

    This architecture is validated through a fully functional prototype,
    verifying its correctness against current Alchemist engine. Evaluation
    demonstrates that while the reactive approach improves modularity and
    expressiveness, it introduces performance compromises. Benchmarking reveals
    that object allocation and notification overheads currently result in lower
    execution speeds compared to the highly optimised original engine,
    particularly in dense dependency graphs. Despite these limitations, this
    work establishes a proof of concept demonstrating that discrete event
    simulation can be effectively modelled through reactive paradigms,
    outlining clear directions for future optimisation.
\end{abstract}
